\section{Theory Bridge and Promotion Discipline}
\label{sec:prof-theory-bridge}

This section distills the formal spine of ORIUS for a professor-facing IEEE draft.  It
connects the battery proof ladder, the universal kernel, and the cross-domain promotion
logic without reintroducing the full monograph narrative.  The key point is that ORIUS
is not presented as a battery theorem generalized by analogy.  It is presented as a
runtime safety grammar whose transfer obligations are explicit, checkable, and separate
from the evidence hierarchy.

\begin{figure}[t]
\centering
\includegraphics[width=\linewidth]{fig_orius_theory_runtime_domain_flow}
\caption{The theory bridge separates three levels of the claim: formal runtime objects,
domain-specific evidence obligations, and promotion status under the defended 3-domain discipline.}
\label{fig:prof-theory-runtime-domain}
\end{figure}

\subsection{Formal setting}
We consider a partially observed safety-constrained control system on a discrete
horizon \(t=0,1,\dots,T-1\):
\[
  x_{t+1}=f(x_t,a_t,\omega_t), \qquad
  y_t=h(x_t,\nu_t), \qquad
  w_t=\psi(y_{0:t}),
\]
with true state \(x_t\in\mathcal{X}\), action \(a_t\in\mathcal{A}\), observation
history \(\mathcal{H}_t=\sigma(y_{0:t})\), and reliability score \(w_t\in[0,1]\).
The safe set is
\[
  \mathcal{S}=\{x\in\mathcal{X}:\phi(x)\ge 0\},
\]
and the episode violation count is
\[
  V_T=\sum_{t=0}^{T-1}\mathbf{1}[x_t\notin\mathcal{S}].
\]
The ORIUS runtime layer takes a nominal candidate action
\(a_t^{\mathrm{cand}}=\pi_t(\mathcal{H}_t)\), constructs an observation-consistent
state set \(U_t\), tightens the admissible action set
\(\mathcal{A}_t^{\mathrm{safe}}(U_t)\), repairs or replaces the candidate action, and
emits a causal certificate.

\subsection{Promotion discipline}
The active ORIUS branch uses a promotion discipline rather than a flat universality gate. A
domain instantiation is not promoted by rhetorical similarity; it is promoted when it supplies:
\begin{enumerate}
  \item a declared observation model and a reproducible replay or experimental protocol;
  \item a sound observation-consistent set \(U_t\) and a domain-specific safe-action set;
  \item a repair or fallback operator that is measurable, causal, and certificate-backed;
  \item a benchmark compatible with the universal schema and a certificate that can be audited;
  \item a domain-specific safety reduction result that can be checked against the shared contract.
\end{enumerate}
This discipline rule is the manuscript's main safeguard: architecture may be universal
before evidence is broad, but the defended empirical claim remains limited to the
promoted three-domain lane until additional domains discharge the same obligations.
For theorem promotion specifically, status is governed by
\texttt{reports/publication/theorem\_promotion\_matrix.csv} and
\texttt{reports/publication/theorem\_result\_cards/*.json}.  The dashboard and
manuscript may call a row flagship only when the statement, assumptions, proof,
code anchor, tests, artifacts, hashes, manuscript anchor, and claim-boundary note
are all present.

\begin{proposition}[Promotion discipline]
\label{prop:universality-gate}
If a domain instantiation satisfies the universal adapter contract and discharges the
promotion obligations above, then the ORIUS kernel can be interpreted on that domain as a
runtime safety layer with the same formal objects used in the battery reference proof:
observation reliability, uncertainty inflation, admissible repair, fallback, and
certificate semantics.  If those obligations are not discharged, the manuscript may
retain only a bounded or gated claim.
\end{proposition}

\begin{proof}[Proof sketch]
The proposition is a packaging statement, not a new control theorem.  It follows directly
from the typed kernel and adapter surface described in the universal-completeness chapter:
the runtime only requires an observation packet, a reliability score, an observation-
consistent state set, a tightened safe-action set, a repair operator, a fallback, and a
causal certificate.  Appendix~A lists these objects and their assumption dependencies.
Once a domain supplies the same objects, the formal language of the battery ladder applies
without changing the runtime contract; what changes is the evidence status of the domain.
\end{proof}

\subsection{T5--T7\_Battery: Temporal certificate and fallback support}
Certificates are temporal objects. T5 (Finite-Horizon Certificate Validity) is a
flagship finite-horizon theorem: if the reachable tube from the
observation-consistent state set stays inside the safe set through horizon~$h$,
the certificate remains valid through that horizon unless invalidating evidence
arrives. T6 (Certificate Expiration Bound) adds a confidence-aware
closed-form first-passage lower-bound proposition on the remaining certificate
life, using the current boundary gap and the drift proxy $\sigma_d$.
T7\_Battery (Feasible Fallback Existence) then closes the operational loop as a
supporting proposition with a piecewise battery fallback: safe hold on the
interior, boundary-aware safe landing near the boundary, and fail-closed behavior
when neither case is certifiable. These surfaces are runtime-enforced:
when the T6 lower bound drops below one step, the guarantee check fails and the shield
must discharge the C7\_BatteryFallback contract.

\subsection{T9: Impossibility of quality-ignorant mandatory release}
T9 is promoted only in its scoped empty-safe-core form.  It is not a universal
claim about controllers allowed to abstain, expand uncertainty, fallback, or deny
release.

The first lower-bound claim is that ignoring reliability is not a benign
simplification when the observation class hides latent states with disjoint safe
actions.

\begin{theorem}[T9: Impossibility of quality-ignorant mandatory release]
\label{thm:prof-t9}
For an observation \(o\), let \(B(o)=\{x:h(x)=o\}\) and
\(K(o)=\cap_{x\in B(o)}C(x)\). If \(B(o)\ne\varnothing\) and
\(K(o)=\varnothing\), no total observation-only mandatory-release controller
can guarantee true-state safety for every \(x\in B(o)\).
\end{theorem}

\begin{proof}[Proof sketch]
The controller must choose one action \(a=\pi(o)\) for the nonempty ambiguity
class. Empty common safe core means no single action belongs to all \(C(x)\) for
\(x\in B(o)\), so some latent state in the ambiguity class makes that mandatory
release unsafe.
\end{proof}

\subsection{T10: Boundary-indistinguishability lower bound}
T10 is a two-state lower-bound theorem.  It does not claim a global minimax
frontier or matching optimality result.

\begin{theorem}[T10: Boundary-indistinguishability lower bound]
\label{thm:prof-t10}
Let \(x_0,x_1\) be boundary states with observation laws \(P_0,P_1\). If
\(d_{\mathrm{TV}}(P_0,P_1)\le\epsilon\) and
\(C(x_0)\cap C(x_1)=\varnothing\), then every observation-only
mandatory-release controller satisfies
\[
  \sup_{i\in\{0,1\}}\Pr_{O\sim P_i}[\pi(O)\notin C(x_i)]
  \ge \frac{1-\epsilon}{2}.
\]
\end{theorem}

\begin{proof}[Proof sketch]
This is a Le Cam style reduction.  For a deterministic release rule, let
\(S_i=\{o:\pi(o)\in C(x_i)\}\).  Disjoint safe-action sets imply
\(S_0\cap S_1=\varnothing\).  A test that decides \(x_0\) on \(S_0\) and
\(x_1\) otherwise succeeds at least as often as the release rule is safe, so the
two-point testing limit \((1+d_{\mathrm{TV}}(P_0,P_1))/2\) upper-bounds release
success.  The displayed lower bound follows by complementing and taking the
larger of the two state risks.  Randomized controllers follow by conditioning
on their internal seed.
\end{proof}

\subsection{Active theorem package after promotion}
The promoted theorem package is deliberately stronger than a relabeling exercise.  Each
surface now has a narrow mathematical statement, an explicit assumption footprint, a
proof source, a code anchor, tests, artifacts, and a claim-boundary note.  The result is
an active register with no defended draft rows.  T5 is a finite-horizon certificate
validity theorem rather than a definition-only row: the certificate is valid exactly
while the reachable tube remains inside the safe set and no invalidating evidence has
arrived.  T8 is a graceful-degradation dominance theorem only after adding useful work:
shutdown may be safe, but it does not satisfy the same safety--work partial order unless
it preserves the required operation fraction.  T9 and T10 are promoted as scoped
lower-bound results over observation-only mandatory release, not as unrestricted
statements about every possible runtime architecture.

The observation-robustness rows follow the same rule.  The Byzantine reliability result
is stated for a trimmed aggregation model with fewer than half corrupted channels; the
corollary to the ORIUS risk envelope is allowed only under that robust-score error
bound.  The stale-observation result is stated as bounded-drift uncertainty growth, not
as a physical law of geometric decay.  The runtime law suite is likewise narrow:
reliability-monotone inflation, safe-set antitonicity, intervention thresholding, and
the observation-ambiguity sandwich are small lemmas and propositions that support the
larger proof spine.  Tminimax is explicitly finite and scoped to two-state boundary
ambiguity classes, while the sensor-necessity row reduces missing critical coordinates
to the empty-safe-core condition.  These constraints are the reason the theorem package
can be ambitious without becoming overbroad.

\subsection{T11: Structural transfer through the universal adapter contract}
The third claim is the one that matters most for a universal paper draft.  It says that
the battery proof pattern can transfer if, and only if, the domain exposes the right
runtime objects.

\begin{theorem}[T11: Structural transfer]
\label{thm:prof-t11}
Let \(\mathcal{D}\) be any domain instantiation that satisfies the universal adapter
contract: it provides an observation packet, a reliability score \(w_t\), an
observation-consistent state set \(U_t\), a domain-safe action set
\(\mathcal{A}_t^{\mathrm{safe}}(U_t)\), a repair map, and a fallback action.  If
\(\Pr[x_t\in U_t\mid\mathcal{H}_t]\ge 1-\alpha\) and every repaired action lies in the
tightened safe-action set, then
\[
  \Pr[x_{t+1}\in\mathcal{S}\mid\mathcal{H}_t]\ge 1-\alpha.
\]
\end{theorem}

\begin{proof}[Proof sketch]
Condition on the observation history.  On the event \(x_t\in U_t\), soundness of the
tightened safe-action set implies that every admissible repaired action preserves the
safe set at the next step.  If the tightened set is empty, the fallback action preserves
safety by contract.  Therefore the only remaining source of failure is the probability
that the true state lies outside \(U_t\), which is at most \(\alpha\).  Appendix~A and
the theorem-surface register record the battery proof instances that instantiate this
logic in the reference domain.
\end{proof}

\begin{corollary}[Episode aggregation under a domain risk budget]
\label{cor:prof-episode-aggregation}
If a domain supplies a predictable per-step risk budget \(r_t\) such that
\(\Pr[x_{t+1}\notin\mathcal{S}\mid\mathcal{H}_t]\le r_t\), then
\[
  \mathbb{E}[V_T]\le \sum_{t=1}^T r_t.
\]
When the budget has the battery-style envelope \(r_t\le \alpha(1-w_t)\), the familiar
upper bound \(\mathbb{E}[V_T]\le \alpha(1-\bar{w})T\) follows immediately.
\end{corollary}

\subsection{Interpretation for the IEEE draft}
The practical message for the IEEE manuscript is that the theory bridge does not need to
be a second thesis chapter.  It needs to do three things well: state the universal
objects, state the universality gate, and state the proof skeleton that connects the
reference battery ladder to any new domain that clears the adapter contract.  That is the
strongest possible version of the claim without inventing unsupported empirical parity.

This is why the theory bridge belongs in the main paper instead of disappearing into a
supplement. The ORIUS claim is not merely that a benchmark improved. It is that
degraded observation changes the structure of what safe release means, and that the
runtime layer makes that change explicit in a way that can be proved, benchmarked, and
audited.
